\section{Integral multiplicities from finite covers}
\label{sec:marks}

In this section take the two-spin germ of Proposition~\ref{prop:spin-sign},
with $w$ the difference of its spin structures and $\epsilon$ its core sign.
Let $\Gamma=\pi_1(C)$ and let $\sigma(H)$ denote the normal determinant
sign on the connected cover corresponding to a finite-index subgroup
$H\le\Gamma$. Write $x$ for a simple nonseparating loop representing
the Poincar\'e dual of $w$ when $w\ne0$, and take $x=1$ when $w=0$.
The covering sign is
\begin{equation}
 \sigma(H)=\epsilon^{[\Gamma:H]}
      \operatorname{sgn}(x\mid\Gamma/H).
 \label{eq:cover-mark}
\end{equation}
For a finite group $Q$, the Burnside ring $A(Q)$ is the Grothendieck
ring of finite $Q$-sets. Its mark $m_H$ at a subgroup $H\le Q$ counts
the $H$-fixed points. Marks on virtual sets take integer values.
For the profinite completion $\widehat\Gamma$ we use the completed
Burnside ring
\[
  \widehat A(\Gamma)=\varprojlim_{N} A(\Gamma/N),
\]
where $N$ ranges over finite-index normal subgroups of $\Gamma$.
The map associated with a quotient $Q'\to Q$ takes fixed points of
its kernel. A transitive orbit survives precisely when its stabilizer
contains the kernel. Thus an element has an integer coefficient
$b_{[H]}$ for each conjugacy class of finite-index subgroups. It can be
read in any quotient through which the action on $\Gamma/H$ factors.
This is the Burnside ring of almost finite sets of
Dress and Siebeneicher~\cite{dress}.

Let $c_m$ denote the signed connected count of \textup{\'e}tale covers
in degree $m$, with the usual inverse deck orders. Let $C_m(h)$ be
the same count with all signs positive for a base of genus $h$.
Recall $h_d=1+d(g-1)$ and the exponents $n_d$ in~\eqref{eq:triangular}.

\begin{theorem}\label{thm:marks}
The signs~\eqref{eq:cover-mark} are the marks of a unit in
$\widehat A(\Gamma)$. Its orbit coefficients satisfy
\begin{equation}
       n_d=\sum_{\substack{[H]\\{[\Gamma:H]=d}}} b_{[H]} .
\label{eq:burnside-primitive}
\end{equation}
Suppose $w\ne0$ and $p$ is an odd prime. Let $M_p$ be the number
of connected degree-$p$ covers with a specified nonzero permutation
sign character in $H^1(C;\mathbb F_2)$. This number is independent
of the chosen nonzero character, and
\begin{equation}
 n_p=-\epsilon\,2^{2g}M_p,\qquad
 |n_p|\ge 2p!\,(p! k_p)^{g-2},
 \label{eq:prime-mark}
\end{equation}
where $k_p$ is the number of partitions of $p$. For $w=0$ one has
$n_p=0$ for every odd prime $p$.
\end{theorem}

\begin{proof}
Fix a quotient $Q$ of $\Gamma$ and denote by $T$ the subgroup
generated by the image of $x$. Write $\epsilon=(-1)^a$ with
$a\in\{0,1\}$. Consider the real permutation representation
\[
       V_Q=\RR[Q/T]\oplus\RR[Q]^{\oplus(1-a)}.
\]
The reduced equivariant Euler characteristic of the one-point
compactification $S^{V_Q}$ is an element $u_Q\in A(Q)$ with mark
$(-1)^{\dim V_Q^H}$ at $H$. Indeed, after equivariant triangulation
and subdivision, each fixed-set Euler characteristic is the
alternating simplex count on the corresponding fixed sphere.
Subtracting the fixed point at infinity gives the reduced Euler
characteristic. This construction of units from real representations
is classical~\cite[Proposition~5.8.2]{tomdieck}.

If $d=[Q:H]$, then $\dim\RR[Q]^H=d$, and
$\dim\RR[Q/T]^H=|H\backslash Q/T|$ is the number $k$ of cycles
of $x$ on $Q/H$. The sign of this permutation is $(-1)^{d-k}$.
Thus the marks of $u_Q$ are precisely~\eqref{eq:cover-mark}.
Marks distinguish elements of $A(Q)$ by triangularity on transitive
orbits. In particular $u_Q^2=1$. Under a further quotient the
mark of the kernel-fixed set at a subgroup downstairs is the old
mark at its inverse image. Formula~\eqref{eq:cover-mark} is identical
on the two coset sets, so the $u_Q$ give the asserted unit.

For completeness let $\mu$ denote the M\"obius function on the
subgroup interval $[H,\Gamma]$. This interval is finite, since it
can be computed in the quotient by the normal core of $H$.
Inverting the marks gives
\begin{equation}
 b_{[H]}=\frac{1}{|N_\Gamma(H)/H|}
                 \sum_{H\le K\le\Gamma}\mu(H,K)\,\sigma(K).
\label{eq:stabilizer-inversion}
\end{equation}
The sum is over subgroups themselves. The numerator counts virtual
points of stabilizer equal to $H$. A single orbit of type $\Gamma/H$
contains $|N_\Gamma(H)/H|$ such points.

The number of subgroups of index $m$ in the genus-$h$ surface group
is $mC_m(h)$, by conjugacy and the deck order. Hence $mc_m$ is the
sum of $\sigma(K)$ over all subgroups $K$ of index $m$ in $\Gamma$.
One orbit $\Gamma/H$ of size $d$ contributes one for each pair
$(K,z)$ with $z\in(\Gamma/H)^K$. For each of the $d$ choices of $z$,
its stabilizer is a surface group of genus $h_d$. The subgroup
$K$ must have index $m/d$ therein. It follows that
\[
 c_m=\sum_{d\mid m}
       \left(\sum_{\substack{[H]\\{[\Gamma:H]=d}}} b_{[H]}\right)
                C_{m/d}(h_d).
\]
All sums are finite since a finitely generated group has finitely
many subgroups of each finite index. Comparison with
\eqref{eq:triangular} gives~\eqref{eq:burnside-primitive} by induction,
using $C_1(h)=1$.

Now $[\Gamma:H]=p$ implies that $H$ is maximal. Its deck group
acts freely on a fiber of $p$ points. If the group is nontrivial,
it acts as a regular cyclic group of order $p$. The monodromy then
lies in the centralizer of a $p$-cycle, all of whose permutations
are even. Consequently any cover on which $x$ acts oddly has trivial
deck group. Formula~\eqref{eq:stabilizer-inversion} gives
$b_{[H]}=-2\epsilon$ for each such cover and zero for the others.
In particular all coefficients vanish in degree $p$ when $w=0$.

If $w\ne0$, the mapping class group acts transitively on nonzero
mod-two cohomology classes and carries covers to covers with the
transformed sign character. Hence each specified nonzero character
occurs on the same number $M_p$ of covers. Exactly $2^{2g-1}$
characters evaluate to one on $x$, which gives the equality
in~\eqref{eq:prime-mark}.

To prove the lower bound, take standard surface generators $a_i,b_i$
with $a_1=x$. Send $a_1$ to any transposition and $a_2$ to any
$p$-cycle in $S_p$. Their images generate $S_p$. Indeed the translates
of the transposition under the cycle are the edges of a connected
graph on $\ZZ/p$, and the edge transpositions of a connected graph
generate the symmetric group. Send each of $b_1,b_2$ to an arbitrary
element centralizing the image of the corresponding $a_i$. Each of
these two handles gives $p!$ choices by the conjugacy class formula.
On each of the remaining $g-2$ handles choose any commuting pair.
There are $p!k_p$ such pairs. The surface relation holds, and every
resulting monodromy is surjective. Relabelling acts freely on these
homomorphisms since the center of $S_p$ is trivial. Dividing their
number by $p!$ gives at least $p!(p!k_p)^{g-2}$ covers on which $x$
acts oddly. Each contributes $-2\epsilon$.
\end{proof}

\begin{example}[Degree three]
For a nonzero character $\alpha$ consider the one-dimensional
$\Gamma$-module $\mathbb F_{3,\alpha}$ on which $s\in\Gamma$ acts
by $(-1)^{\alpha(s)}$. Identify $S_3$ with the affine group of
$\mathbb F_3$. Specifying monodromy with sign $\alpha$ amounts to
specifying its translation part, a cocycle in
$Z^1(\Gamma;\mathbb F_{3,\alpha})$. The image fixes a point if and
only if the cocycle is a coboundary. Every intransitive action on
three letters has a fixed point. Conjugation by translations changes
a cocycle by a coboundary, and conjugation by reflection negates its
class. The nonzero classes therefore give transitive covers in pairs.
The local system has no invariants and is self-dual, so $H^0$ and
$H^2$ vanish and $\dim H^1=2g-2$ by the Euler characteristic.
We obtain, for $w\ne0$,
\[
 M_3=\frac{3^{2g-2}-1}{2},\qquad
 n_3=-\epsilon\,2^{2g-1}(3^{2g-2}-1).
\]
In genus two the exponent is $-64\epsilon$.
\end{example}

\begin{remark}[Comparison with Parker's signs]
Parker records determinant signs on all subgroup-fixed subspaces of
the deck group in the integral algebra on subgroups whose product
is intersection~\cite[Definitions~1.1--1.2 and Lemma~1.3]{parker}.
The map from $A(Q)$ to that algebra sends a finite $Q$-set $X$ to
$\sum_{z\in X}[Q_z]$, where $Q_z$ is its point stabilizer.
Intersection computes the stabilizer in a product, and the marks
recover the virtual set. Under this injection $u_Q$ gives Parker's
sign element on the regular $Q$-cover, since both give the same
determinant on every fixed subspace. Parker's local contribution of
an individual cover uses its own deck group. For instance a cover
of odd prime degree on which $x$ acts oddly has trivial deck group
and local sign $-\epsilon$. Its orbit coefficient in the regular
closure in~\eqref{eq:stabilizer-inversion} is $-2\epsilon$.
\end{remark}

\begin{corollary}\label{cor:general-sign}
Let $C$ be any isolated genus-$g$ curve satisfying
Proposition~\ref{prop:factors}, and let $\epsilon$ be its normal
determinant sign. At every odd prime $p$ its elementary exponent
at $(p,1+p(g-1))$ satisfies
\[
 n_p=-2\epsilon\,\#\{[f]:\deg f=p,\ f\text{ connected \'etale},
                      \ \sigma_D(f)=-\epsilon\}.
\]
In particular $-\epsilon n_p$ is an even nonnegative integer.
Every connected cover with odd-order monodromy has sign $\epsilon$.
\end{corollary}

\begin{proof}
For the absolute signs use the coincident-spin operator of
Proposition~\ref{prop:spin-sign} as reference. Its two identical
determinant signs multiply to $+1$ on every unramified cover.
The relative unit from Theorem~\ref{thm:endpoints} with this initial
operator therefore has the prescribed absolute marks.

For $p$ prime, \eqref{eq:triangular} subtracts the unsigned degree-$p$
count times $\epsilon$ from the signed count. An index-$p$ cover
with deck transformations is regular cyclic, as in the proof of
Theorem~\ref{thm:marks}. On a regular cyclic $p$-cover the cardinality
of a finite $C_p$-set is congruent modulo $p$ to its number of fixed
points. This holds for virtual sets as well. Since $p$ is odd, the
two signs of the unit $u_{C_p}(D)$ agree. Every surviving term in the
difference thus has trivial deck group, proving the formula.

More generally, on an odd-order monodromy group the trivial representation
is the only irreducible representation of real type. Indeed squaring
permutes the elements of such a group, so the Frobenius--Schur indicator
of each nontrivial irreducible complex character vanishes by character
orthogonality. In the construction of Theorem~\ref{thm:endpoints} only
the trivial representation can contribute a sign, and its fixed space
has dimension one at every subgroup. The sign is therefore the base sign.
\end{proof}

\subsection*{The full product on an isolating interval}
There is a separate all-genus conclusion when a neighborhood isolates
the entire family. We state it with the compactness hypothesis
required for the interval relation. Say that a normalized series
$R(q)=1+\sum_{m>0}R_mq^m$ in $H[[q]]$ respects the cycle operations
if addition and multiplication of cycles induce
\[
 \Delta R_m=\sum_{i+j=m}R_i\otimes R_j,\qquad
 \rho_k R_m=\begin{cases}R_{m/k}&k\mid m,\\0&k\nmid m.\end{cases}
\]
Balanced products and their integral powers have these properties
\cite[\S7]{WueProduct}. Substitution $q\mapsto q^r$ preserves them:
for the division operation on the substituted series this follows
directly by asking whether $kr$ divides the original label $m$.
The coefficient at label $m$ on either side is zero in the other
cases.

\begin{proposition}\label{prop:full-transport}
The balanced lift in Proposition~\ref{prop:factors} is the unique
normalized lift respecting the cycle operations. Suppose an interval
of tame almost-complex neighborhoods carries an additive positive
integral degree label. At each end let the selected connected supports
be finitely many pairwise disjoint curves satisfying
Proposition~\ref{prop:factors}, of positive labels $r_i$ at the left
end and $s_j$ at the right end. Assume, at each finite degree cutoff,
that the characteristic insertion of the cycles in the selected
neighborhood is compact open and closed over the interval; it selects
all sums of multiples of these supports at the endpoints, up to the
cutoff. Then in $H[[q]]$ one has the full equality
\[
        \prod_i \widetilde P_i(q^{r_i})
        =\prod_j \widetilde P_j(q^{s_j}).
\]
\end{proposition}

\begin{proof}
For two lifts with the given operations, a first nonzero difference
is primitive by the coproduct formula and is killed by all
$\rho_k$ for $k>1$. It is a finite rational linear combination
of $x_{h,1}$ by \cite[Lemma~7.1]{WueProduct}. Its image in $S$ vanishes,
which contradicts \cite[Theorem~5.1]{WueProduct}.
For the interval the family insertion transports the entire selected
series in $S$ at each cutoff. At the ends disjoint open excision
factors the insertion into the product of the local insertions,
where the substitutions retain the positive labels. Apply the
first-difference argument to the two lifted products in each
degree. The same proof works after concatenation, or on either
boundary path of a compactly isolated two-parameter family.
\end{proof}

One sufficient form of the interval hypothesis is a closed tame
ambient family with a relatively compact spatial tube and a
fixed energy bound for each degree cutoff, in which no selected cycle
can reach the tube boundary while remaining in its closure.
Bounded-energy cycle compactness and Hausdorff convergence give
the required compact insertion as in Proposition~\ref{prop:factors}.
The fixed-cover bifurcation theorem gives no such all-genus spatial
isolation along its full path. Indeed, in its linear realization a
nonzero solution can be scaled all the way to the boundary.

\subsection*{A universal correction on the specified cover}
For a crossing of Theorem~\ref{thm:walls} fix $f:Y_m\to C$ and
pull back the tube to $Y_m$. In this tube the cores of the graph
equation over $f$ become embedded sections. The unit inclusion
of an isolated such section of degree one defines an insertion even
when its higher multiples have not been isolated. Write
$x_{h,1}$ for the degree-one genus-$h$ class as in
Section~\ref{sec:factors}.

\begin{corollary}\label{cor:marked-wall}
In the pulled-back tube, the change of the sum of the degree-one
classes of the nonzero graphs at a crossing is
$(\sigma(f)^--\sigma(f)^+)x_{h_m,1}$ in $H$.
Counting these graphs modulo reparametrization over $C$ divides this
change by $|\operatorname{Aut}_C(f)|$.
\end{corollary}

\begin{proof}
A compact effective cycle of projection degree one in the pulled-back
tube has precisely one irreducible component, of multiplicity one.
A vertical component in an affine fiber is impossible, and the normalized
degree-one projection is an isomorphism. Hence the entire degree-one
cycle is the graph of a smooth section of $Y_m$. Near the crossing
the finite-order reduction in Section~\ref{sec:walls} isolates these
sections. More explicitly, bounded sections whose sup norms tend to
zero satisfy the fixed-cover semilinear equation with a right hand
side tending uniformly to zero. The elliptic estimate for the fixed
Dolbeault operator then gives convergence to zero in $W^{1,r}$ for
$r>2$. This places them in the reduction chart. Conversely its
solutions converge smoothly, by elliptic regularity. A small tube
about zero and then a sufficiently short interval isolate the
selected graphs in degree one. At the central parameter there is
no other small solution, and compactness of bounded sections follows
by the same estimate and elliptic bootstrapping.

For an isolated regular embedded graph of genus $h$ its degree-one
insertion equals $\operatorname{sign}(D)\Phi(x_{h,1})$ in $S$.
Indeed work at the energy of the primitive core in the compact
realization of Proposition~\ref{prop:factors}. At the original
endpoint the degree-one core is unique in an isolating tube and is
regular. One can choose the nearby generic structure fixing it as in
\cite[proof of Proposition~3.54]{diw}. In the short transport it is
the unique nearby curve of its class, by the openness of the embedded multiplicity-one locus
\cite[Proposition~2.21]{diw} and the parameterized implicit function theorem.
The cluster isotopy to the elementary endpoint has only corrections
in proper multiples, all above this cutoff. It transports the leading
factor with the orientation sign of the core.
The unique lift of this degree-one class compatible with the operations
is the same signed $x_{h,1}$, by the primitive argument of
Proposition~\ref{prop:full-transport} at degree one.
Apply this fact to each regular graph on the two adjacent slices
and the signed count of Theorem~\ref{thm:walls}. Each nonzero simple
support with projection $f$ has exactly $|\operatorname{Aut}_C(f)|$
parametrizations in the fixed-cover chart. This proves both statements.
\end{proof}

This equality retains the universal class of the embedded graph on the
specified cover. Recovering the entire degree-$m$ coefficient of the
original, unmarked tube would also require the other projection types.
